\section{大规模训练的工程挑战}

\subsection{硬件故障与容错}

在 24,000 块 GPU 上训练数月，硬件故障是必然的：

\begin{figure}[H]
\centering
\includegraphics[width=0.95\textwidth]{slides-images/slide-046.jpg}
\caption{Llama 3 训练期间的故障统计：54 天内发生 466 次任务中断，其中 78\% 可归因于已知原因\protect\footnotemark}
\end{figure}
\footnotetext{Slides 第46页。}

\begin{warningbox}{大规模训练的故障率}
Llama 3 在 54 天训练中经历了 \textbf{466 次任务中断}，平均每天 8--9 次。硬件故障占 58.7\%（GPU 故障 30.1\%、主机维护 14.1\%、网络问题 5.5\%）。这意味着训练系统必须能够\textbf{自动检测故障、重启并从检查点恢复}，否则大规模训练根本不可能完成。
\end{warningbox}

\subsection{本章小结}

大规模训练不仅是算法和硬件的挑战，更是系统工程的挑战。自动故障检测、快速检查点恢复和训练监控是使训练在数万 GPU 上稳定运行数月的关键基础设施。

%% ========== 总结与延伸 ==========

